% =====================================================================
%  Chapter 3 -- Materials and Method
%  Status: SKELETON (headings only)  --  see status.md
%  Sources: Capstone B report Ch. 4-6; master_report_metadata/
%           (tables/table03, table04; A_generated_figures/02_methodology,
%           03_training_curves); notebooks/*; documents/project_README.md
% =====================================================================
\chapter{Materials and Method}
\label{ch:method}

\todo{Chapter introduction.}

% =====================================================================
\section{Materials}
\label{sec:materials}

\subsection{Dataset Acquisition}
\label{sec:dataset-acquisition}
% ASL Citizen download + licence (Microsoft Research License Terms,
% no redistribution); why WLASL was abandoned (median 6 videos/class).
\todo{Write.}

\subsection{Dataset Preprocessing}
\label{sec:preprocessing}
% Manifest from official train/val/test CSVs; 63,550 cached videos /
% 2,731 glosses; 3 cache shards on Kaggle; top-K selection; class-list
% fingerprint 72d3ef78291c shared by all 1000-class models; 16 frames.
\todo{Write.}

\subsection{Feature Extraction}
\label{sec:feature-extraction}
% MediaPipe Tasks (Hand/Pose/Face landmarkers, VIDEO mode); 58 points x
% (x,y) = 116 features; shoulder-centre / shoulder-width normalisation
% (Eq.); velocity; 112x112 signer crop with 25% margin; x-position hand
% assignment; detected-frame sampling history (38.1% -> 30.7%).
% Figure: B_original_notebook_figures/... hand landmark schematic.
\todo{Write.}

% =====================================================================
\section{Methods}
\label{sec:methods}

\subsection{Why Transformer Models instead of RNN or Encoder--Decoder}
\label{sec:why-transformer}
\todo{Write.}

\subsection{Proposed Models}
\label{sec:proposed-models}
% Development path: Capstone B iterations (summarise the B report Ch. 5
% by WHAT changed and WHY, never by version number) and the post-B
% redesign: geometry-only ablation, ST-GCN, Video Swin-T.
% Do NOT write internal labels such as "v9", "v10" or notebook names.
\todo{Write.}

\subsection{Model Architectures}
\label{sec:architectures}
% Figure: A_generated_figures/02_methodology/fig02_02_architecture_overview
% Four architectures: R3D-18 Dual-Stream, Geometry Transformer, ST-GCN,
% Video Swin-T (see metadata/project_metadata.json "models").
\todo{Write.}

\subsection{Algorithm / Model Formulation: Attention Models}
\label{sec:attention}
\todo{Write (scaled dot-product attention, multi-head attention,
positional encoding, shifted-window attention, graph convolution).}

\subsection{Transformers}
\label{sec:transformers}
\todo{Write.}

\subsection{Proposed Model Design / Framework}
\label{sec:framework}
% Figure: A_generated_figures/02_methodology/fig02_01_recognition_pipeline
% (already used in Ch.1? -> NO, Ch.1 uses the app screenshot; use it here)
\todo{Write.}

\subsection{Experimental Setup}
\label{sec:experimental-setup}
\todo{Write.}

\subsection{Evaluation Metrics}
\label{sec:metrics}
% top-1, top-5, macro precision/recall/F1, weighted F1, CPU latency.
\todo{Write.}

\subsection{Hardware, Design Tools, and Development Environment}
\label{sec:hardware}
% Kaggle T4 GPUs (Swin: 2x T4, 6.1 h); local Windows 11 CPU; Python
% 3.13.5, torch 2.15.0+cpu, mediapipe 0.10.35, streamlit 1.67.0 ...
\todo{Write.}

\subsection{Dataloader, Optimizer, and Loss Function}
\label{sec:optim}
\todo{Write.}

\subsection{Model Parameters}
\label{sec:model-params}
% tables/table03_training_hyperparameters.tex
\todo{Write.}

\subsection{Computation of Parameters}
\label{sec:param-computation}
\todo{Write.}

% =====================================================================
\section{Designing and Building the Web Application}
\label{sec:webapp-design}
% Streamlit pages: Home, Recognize, Results, Methodology, About.
% Screenshots: C_app_screenshots/app01..app12.
\todo{Write.}

\section{Integrating the Models with the Web Application}
\label{sec:webapp-integration}
% video_processor.py (port of the cache builder), model_loader.py
% (strict load, HF download, one model in memory), predictor.py.
\todo{Write.}

\section{Summary}
\label{sec:method-summary}
\todo{Write.}
